% from main.tex
%%%%%%%%%%  scan idx 124  |  folio 111  %%%%%%%%%%
le
foncteur dérivé du produit tensoriel \((E',E)\longmapsto E'\otimes_{f^{-1}(A)} f^{-1}(E)\). Soient \(X\in\Ob S\), \(X'=u(X)\), \(E'\in\Ob\Dpx^-(_{A'}X'_A)\), \(E\in\Ob\Dpx^-(_A X)\). Alors si \(E'\) est \(m'\)-pseudo-cohérent (comme objet de \(\Dpx^-(_A S')\)) et acyclique en degré \(>n'\) et si \(E\) est \(m\)-pseudo-cohérent et acyclique en degré \(>n\), \(E'\otimes_A^{\mathrm L}E\) est \(\sup(m+n',m'+n)\)-pseudo-cohérent et acyclique en degré \(>n+n'\).

\emph{Preuve.} Quitte à remplacer \(S\) par \(S_{/X}\), on peut supposer \(X=S\), et l'on applique (2.13) avec \(u=E'\otimes_A^{\mathrm L}:\Dpx^-(_A S)\to\Dpx^-(_{A'}S')\). Soient \(U\in\Ob S\), \(U'=f^{-1}(U)\), \(E\) un complexe de \(A_U\)-Module à gauche. Si \(E\) est acyclique en degré \(>n\), \(E'\otimes_A^{\mathrm L}E\) est acyclique en degré \(>n+n'\) (définition de \(\otimes_A^{\mathrm L}\)); d'autre part si \(E=(A_U)^p\), alors \(E'\otimes_A^{\mathrm L}E=(E'|U')^p\), qui est \(m'\)-pseudo-cohérent par hypothèse (et (2.12)). Les conditions de (2.13) sont donc bien vérifiées, avec \(a=m'\), \(b=n'\), d'où la conclusion.

En particulier :

\textbf{Corollaire 2.16.1.} Sous les hypothèses de (2.16), on a :

\textbf{a)} Le foncteur \(\mathrm Lf^*: \Dpx^-(_A S)\to \Dpx^-(_{A'}S')\) (qui n'est autre que le foncteur \(A'\otimes_A^{\mathrm L}\)) induit un foncteur
\[
  \mathrm Lf^*: \Dpx^-(_A S)_{\coh}\longrightarrow \Dpx^-(_{A'}S')_{\coh}
\]
(notation de (2.6)).

\textbf{b)} Si \(A\) est commutatif, le foncteur \(\otimes_A^{\mathrm L}:\Dpx^-(S)\times_S\Dpx^-(S)\longrightarrow \Dpx^-(S)\) induit un foncteur
\[
  \otimes_A^{\mathrm L}:\Dpx^-(S)_{\coh}\times_S\Dpx^-(S)_{\coh}
  \longrightarrow \Dpx^-(S)_{\coh} .
\]

